\documentclass[10pt,a4paper]{amsart}
\usepackage[margin=19mm]{geometry}
\usepackage{amsmath,amssymb,mathtools}
\usepackage[scr=rsfs,cal=euler]{mathalfa}
\usepackage{xfrac}
\usepackage[unicode,hidelinks,bookmarks=false]{hyperref}
\newcommand{\ie}{i.e.\ }
\newcommand{\cf}{cf.\ }
\newcommand{\resp}{respectively\ }
\newcommand{\Subsection}{\subsection}
\providecommand{\nobreakdash}{\nobreak\hbox{-}}
\setlength{\emergencystretch}{2em}
\multlinegap0pt
\usepackage{amssymb}
\usepackage{tikz}
\newtheorem{proposition}{Proposition}
\newcommand{\supp}{\operatorname{supp}}
\title{Sharp Ill-Posedness of the Euler Equations in Lorentz Spaces}
\author{Jeaheang Bang, Alexey Cheskidov}
\date{Source: arXiv:2605.16502v1 (CC BY 4.0)}
\begin{document}
\maketitle

\noindent\emph{Source and licence.} Sharp Ill-Posedness of the Euler Equations in Lorentz Spaces. Version arXiv:2605.16502v1 (CC BY 4.0), distributed by arXiv under the Creative Commons Attribution 4.0 International licence, http://creativecommons.org/licenses/by/4.0/, as recorded in the arXiv licence metadata for this exact version and shown on the abstract page https://arxiv.org/abs/2605.16502v1. Full licence notice: \texttt{licenses/arxiv-2605.16502-CC-BY-4.0.txt}.\par\smallskip

\noindent\textit{Editorial scope.} Counted unit: the article's proposition prop:transport\_by\_flow (source0.tex lines 7164--7233). The article's complete original proof of it is delivered with it (source0.tex lines 7235--7625, one \texttt{proof} environment of 391 lines). No mathematics is written, repaired or re-derived: the statement and the proof are the article's own lines, in the article's own notation, and no retained line is substituted or re-flowed.  No further source-local context is needed: every cross-reference of every retained line resolves inside the two delivered spans.  Nothing was omitted from any retained span, so the package carries no omission record.  No retained line contains a citation, so the wrapper carries no bibliography and no bibliography entry.  No retained line contains a figure environment, an image-inclusion command or drawing code, so no image is delivered and the package declares no asset dependency.  Editorial formatting only, in addition: an AMS \texttt{amsart} preamble that restates the article's own declarations and macros used by the retained lines, namely the environments \texttt{proposition} on separate counters, together with the standard packages those lines need.  The retained environments print in this document as Proposition 1 in order of appearance, whereas the article prints the same items with the numbering of its own full document.  The complete native source of the article as distributed in the arXiv e-print for this exact version is bundled unchanged as \texttt{sources/200718/source0.tex}.
\begin{proposition}[Cauchy formula of the transport equation along a given flow]
\label{prop:transport_by_flow}
Let $T_0>0$. Let
\[
u\in L^\infty\bigl((0,T_0)\times\mathbb R^3\bigr)
\]
be divergence-free, and assume that there exists a flow map
\[
Y:\mathbb R^3\times[0,T_0]\to\mathbb R^3
\]
such that:

\smallskip
\noindent
(i) for every $a\in\mathbb R^3$, the map $t\mapsto Y(a,t)$ is absolutely continuous and
\[
\partial_tY(a,t)=u\big(Y(a,t),t\big)
\qquad\text{for a.e. }t\in(0,T_0),
\qquad
Y(a,0)=a;
\]

\smallskip
\noindent
(ii) for every $t\in[0,T_0]$, the map $Y_t:=Y(\cdot,t)$ is a measure-preserving
homeomorphism of $\mathbb R^3$ onto itself;

\smallskip
\noindent
(iii) there exists $\gamma\in(0,1]$ such that for every compact set $K\subset\mathbb R^3$,
\[
\sup_{t\in[0,T_0]}
\Bigl(
[Y_t]_{C^\gamma(K)}+[Y_t^{-1}]_{C^\gamma(K)}
\Bigr)
<\infty .
\]

Let $\xi_0\in L^1(\mathbb R^3)$, and let
\[
\xi\in L^1_{\mathrm{loc}}\bigl([0,T_0)\times\mathbb R^3\bigr),
\qquad
u\xi\in L^1_{\mathrm{loc}}\bigl([0,T_0)\times\mathbb R^3\bigr)
\]
satisfy
\begin{equation}
\label{eq:transport_distributional_appendix}
\int_0^{T_0}\int_{\mathbb R^3}
\xi(x,t)\bigl(\partial_t\psi(x,t)+u(x,t)\cdot\nabla\psi(x,t)\bigr)\,dx\,dt
+\int_{\mathbb R^3}\xi_0(x)\psi(x,0)\,dx
=0
\end{equation}
for every $\psi\in C_c^\infty(\mathbb R^3\times[0,T_0))$.

Assume in addition that there exists a constant $C_u>0$ such that
\begin{equation}
\label{eq:u_loglip_appendix}
|u(x,t)-u(y,t)|
\le
C_u |x-y|\log\!\left(2+\frac{1}{|x-y|}\right)
\end{equation}
for all $x,y\in\mathbb R^3$ and for a.e. $t\in(0,T_0)$.

Then
\[
\xi(x,t)=\xi_0\bigl(Y_t^{-1}(x)\bigr)
\qquad
\text{for a.e. }(x,t)\in\mathbb R^3\times(0,T_0).
\]
\end{proposition}

\begin{proof}
Define
\[
\widetilde{\xi}(x,t):=\xi_0\bigl(Y_t^{-1}(x)\bigr).
\]
Since each $Y_t$ is measure-preserving,
\[
\|\widetilde{\xi}(t)\|_{L^1(\mathbb R^3)}=\|\xi_0\|_{L^1(\mathbb R^3)}
\qquad\text{for every }t\in[0,T_0],
\]
hence
\[
\widetilde{\xi}\in L^\infty\bigl(0,T_0;L^1(\mathbb R^3)\bigr)
\subset L^1_{\mathrm{loc}}\bigl((0,T_0)\times\mathbb R^3\bigr).
\]

Fix $\eta\in C_c^\infty(\mathbb R^3)$ and $\chi\in C_c^\infty(0,T_0)$. Set
\[
K:=\operatorname{supp}\eta,
\qquad
M:=\|u\|_{L^\infty((0,T_0)\times\mathbb R^3)}.
\]
Define
\[
K_*:=\{x\in\mathbb R^3:\operatorname{dist}(x,K)\le MT_0\},
\qquad
K_{**}:=\{x\in\mathbb R^3:\operatorname{dist}(x,K)\le 2MT_0\}.
\]
For $(x,t)\in\mathbb R^3\times[0,T_0]$, define
\[
\Phi(x,t):=
\int_t^{T_0}\chi(\tau)\,\eta\Bigl(Y_\tau \big(Y_t^{-1}(x)\big)\Bigr)\,d\tau.
\]

We first record some elementary properties of $\Phi$.

\smallskip
\noindent\emph{Step 1: support and H\"older regularity of $\Phi$.}
If $\Phi(x,t)\neq 0$, then there exists $\tau\in \operatorname{supp}\chi$ such that
$Y_\tau(Y_t^{-1}(x))\in K$. Writing $a:=Y_t^{-1}(x)$ and using the flow equation,
\[
|x-Y_\tau(a)|=|Y_t(a)-Y_\tau(a)|
\le M|t-\tau|
\le MT_0,
\]
hence $x\in K_*$. Therefore
\[
\operatorname{supp}\Phi(\cdot,t)\subset K_*
\qquad\text{for every }t\in[0,T_0].
\]

Next, if $x\in K_*$ and $a=Y_t^{-1}(x)$, then
\[
\operatorname{dist}(a,K)\le \operatorname{dist}(a,x)+\operatorname{dist}(x,K)\le 2MT_0,
\]
so
\[
Y_t^{-1}(K_*)\subset K_{**}
\qquad\text{for every }t\in[0,T_0].
\]
By assumption (iii), there exists $\beta:=\gamma^2\in(0,1]$ and a constant $C>0$
depending only on $K_*,K_{**}$ such that for all $t,\tau\in[0,T_0]$,
\[
\left|Y_\tau\big(Y_t^{-1}(x_1)\big)-Y_\tau\big(Y_t^{-1}(x_2)\big)
\right|
\le C |x_1-x_2|^\beta,
\qquad x_1,x_2\in K_*.
\]
Hence
\[
\bigl[\eta\circ Y_\tau\circ Y_t^{-1}\bigr]_{C^\beta(K_*)}
\le C\|\nabla\eta\|_{L^\infty},
\]
uniformly in $t,\tau\in[0,T_0]$. It follows that
\begin{equation}
\label{eq:Phi_holder_appendix}
\sup_{t\in[0,T_0]}[\Phi(\cdot,t)]_{C^\beta(K_*)}
\le C\|\chi\|_{L^1(0,T_0)}\|\nabla\eta\|_{L^\infty}.
\end{equation}



Since $\supp \Phi(\cdot,t)\subset K_*$ and $K_*$ is closed, we have
$\Phi(\cdot,t)=0$ on $\mathbb R^3\setminus K_*$. By continuity, it also follows that
$\Phi(\cdot,t)=0$ on $\partial K_*$. Consequently,
\[
[\Phi(\cdot,t)]_{C^\beta(\mathbb R^3)}
\le
[\Phi(\cdot,t)]_{C^\beta(K_*)}.
\]
Indeed, if $x,y\in K_*$ or $x,y\notin K_*$, the claim is immediate. If $x\in K_*$ and
$y\notin K_*$, choose $z\in [x,y]\cap\partial K_*$
 where $[x,y]$ denotes the straight line segment connecting $x$ to $y$.
Then $\Phi(z,t)=\Phi(y,t)=0$, so
\[
|\Phi(x,t)-\Phi(y,t)|
=
|\Phi(x,t)-\Phi(z,t)|
\le
[\Phi(\cdot,t)]_{C^\beta(K_*)}\,|x-z|^\beta
\le
[\Phi(\cdot,t)]_{C^\beta(K_*)}\,|x-y|^\beta.
\]
The case $x\notin K_*$ and $y\in K_*$ is symmetric. Combining this with
\eqref{eq:Phi_holder_appendix}, we obtain
\begin{equation}
\label{eq:Phi_holder_global_appendix}
\sup_{t\in[0,T_0]}[\Phi(\cdot,t)]_{C^\beta(\mathbb R^3)}
\le
C\|\chi\|_{L^1(0,T_0)}\|\nabla\eta\|_{L^\infty}.
\end{equation}


\smallskip
\noindent\emph{Step 2: $\Phi$ solves a backward transport equation.}
Let $\varphi\in C_c^\infty(\mathbb R^3\times(0,T_0))$. By the change of variables
$x=Y_t(a)$ and the measure-preserving property of $Y_t$,
\begin{multline*}
\int_0^{T_0}\int_{\mathbb R^3}
\Phi(x,t)\bigl(\partial_t\varphi(x,t)+u(x,t)\cdot\nabla\varphi(x,t)\bigr)\,dx\,dt 
\\
=
\int_{\mathbb R^3}\int_0^{T_0}
\left(\int_t^{T_0}\chi(\tau)\eta(Y_\tau(a))\,d\tau\right)
\frac{d}{dt}\varphi(Y_t(a),t)\,dt\,da .
\end{multline*}
Set
\[
G_a(t):=\int_t^{T_0}\chi(\tau)\eta(Y_\tau(a))\,d\tau.
\]
Since $G_a$ is absolutely continuous and
\[
G_a'(t)=-\chi(t)\eta(Y_t(a))
\qquad\text{for a.e. }t,
\]
integration by parts in $t$ gives
\begin{align*}
\int_0^{T_0}\int_{\mathbb R^3}
\Phi(x,t)\bigl(\partial_t\varphi(x,t)+u(x,t)\cdot\nabla\varphi(x,t)\bigr)\,dx\,dt 
&=
-\int_{\mathbb R^3}\int_0^{T_0}G_a'(t)\varphi(Y_t(a),t)\,dt\,da \\
&=
\int_{\mathbb R^3}\int_0^{T_0}\chi(t)\eta(Y_t(a))\varphi\big(Y_t(a),t\big)\,dt\,da \\
&=
\int_0^{T_0}\int_{\mathbb R^3}\chi(t)\eta(x)\varphi(x,t)\,dx\,dt .
\end{align*}
Therefore, in distributions on $(0,T_0)\times\mathbb R^3$,
\begin{equation}
\label{eq:Phi_backward_appendix}
\partial_t\Phi+u\cdot\nabla\Phi=-\chi(t)\eta(x).
\end{equation}

\smallskip
\noindent\emph{Step 3: spatial mollification and commutator.}
Let $\rho\in C_c^\infty(\mathbb R^3)$ be a standard nonnegative mollifier supported in
$B_1(0)$ with $\int\rho=1$, and set
\[
\rho_\varepsilon(x):=\varepsilon^{-3}\rho(x/\varepsilon),
\qquad
\Phi_\varepsilon:=\rho_\varepsilon * \Phi,
\qquad
\eta_\varepsilon:=\rho_\varepsilon * \eta.
\]
Then $\Phi_\varepsilon$ is smooth in $x$, compactly supported in $K_*+B_\varepsilon$, and
from \eqref{eq:Phi_backward_appendix},
\[
\partial_t\Phi_\varepsilon+\rho_\varepsilon*(u\cdot\nabla\Phi)
=
-\chi(t)\eta_\varepsilon .
\]
Hence
\begin{equation}
\label{eq:Phi_eps_eq_appendix}
\partial_t\Phi_\varepsilon+u\cdot\nabla\Phi_\varepsilon
=
-\chi(t)\eta_\varepsilon+R_\varepsilon,
\end{equation}
where
\[
R_\varepsilon
:=
u\cdot\nabla\Phi_\varepsilon-\rho_\varepsilon*(u\cdot\nabla\Phi).
\]
Using $\operatorname{div}u=0$, we may rewrite $R_\varepsilon$ as
\[
R_\varepsilon(x,t)
=
\int_{\mathbb R^3}
\nabla\rho_\varepsilon(y)\cdot\bigl(u(x,t)-u(x-y,t)\bigr)
\bigl(\Phi(x-y,t)-\Phi(x,t)\bigr)\,dy .
\]
Indeed, the term containing $\Phi(x,t)$ alone vanishes because
\[
\int_{\mathbb R^3}\nabla\rho_\varepsilon(y)\cdot\bigl(u(x,t)-u(x-y,t)\bigr)\,dy
=
-\int_{\mathbb R^3}\rho_\varepsilon(y)\,\operatorname{div}u(x-y,t)\,dy
=0.
\]



If $x\notin K_*+\overline B_\varepsilon$, then $\operatorname{dist}(x,K_*)>\varepsilon$. Since
$\supp \nabla\rho_\varepsilon\subset B_\varepsilon(0)$, for every
$y\in\supp \nabla\rho_\varepsilon$ we have
\[
\operatorname{dist}(x-y,K_*)\ge \operatorname{dist}(x,K_*)-|y|>0.
\]
Hence $\Phi(x,t)=\Phi(x-y,t)=0$, and therefore $R_\varepsilon(x,t)=0$. Thus
\begin{equation}
\label{eq:R_eps_support_appendix}
\supp R_\varepsilon(\cdot,t)\subset K_*+\overline B_\varepsilon
\qquad\text{for every }t\in[0,T_0].
\end{equation}

Using \eqref{eq:u_loglip_appendix} and \eqref{eq:Phi_holder_global_appendix}, we obtain
for all $(x,t)\in\mathbb R^3\times(0,T_0)$
\begin{align*}
|R_\varepsilon(x,t)|
&\le
C[\Phi(\cdot,t)]_{C^\beta(\mathbb R^3)}
\int_{\mathbb R^3}
|\nabla\rho_\varepsilon(y)|\,|y|^\beta\,|u(x,t)-u(x-y,t)|\,dy \\
&\le
C\int_{|y|<\varepsilon}
|\nabla\rho_\varepsilon(y)|\,|y|^{1+\beta}
\log\!\left(2+\frac1{|y|}\right)\,dy \\
&\le
C\,\varepsilon^\beta\log\!\left(2+\frac1\varepsilon\right).
\end{align*}
Therefore
\begin{equation}
\label{eq:R_eps_to_zero_appendix}
\|R_\varepsilon\|_{L^\infty((0,T_0)\times\mathbb R^3)}
\longrightarrow 0
\qquad\text{as }\varepsilon\to 0.
\end{equation}



\smallskip
\noindent\emph{Step 4: testing the weak formulation with $\Phi_\varepsilon$.}
Since $\chi\in C_c^\infty(0,T_0)$, there exists $\delta_0\in(0,T_0)$ such that
\[
\chi(t)=0
\qquad\text{for }t\in[T_0-\delta_0,T_0].
\]
Hence, by the definition of $\Phi$,
\[
\Phi(x,t)=0
\qquad\text{for all }(x,t)\in\mathbb R^3\times[T_0-\delta_0,T_0],
\]
and therefore the same holds for $\Phi_\varepsilon$.

Also, \eqref{eq:Phi_eps_eq_appendix} yields
\[
\partial_t\Phi_\varepsilon
=
-u\cdot\nabla\Phi_\varepsilon-\chi(t)\eta_\varepsilon+R_\varepsilon
\in L^\infty\bigl((0,T_0)\times\mathbb R^3\bigr),
\]
so $\Phi_\varepsilon$ is compactly supported and belongs to
$W^{1,\infty}(\mathbb R^3\times[0,T_0])$.

Choose $\theta\in C_c^\infty((-1,1))$ such that
\[
\theta\equiv 1
\qquad\text{on }[-\tfrac12,\tfrac12],
\]
and define $\widetilde{\Phi}_\varepsilon$ on $\mathbb R^3\times\mathbb R$ by
\[
\widetilde{\Phi}_\varepsilon(x,t):=
\begin{cases}
\theta(t)\Phi_\varepsilon(x,0), & t<0,\\[1mm]
\Phi_\varepsilon(x,t), & 0\le t\le T_0,\\[1mm]
0, & t>T_0.
\end{cases}
\]
Then $\widetilde{\Phi}_\varepsilon\in W^{1,\infty}(\mathbb R^3\times\mathbb R)$ and has
compact support.

Let $\vartheta\in C_c^\infty(\mathbb R^4)$ be a standard mollifier, and set
\[
\vartheta_\sigma(z):=\sigma^{-4}\vartheta(z/\sigma),
\qquad
\Psi_{\varepsilon,\sigma}:=\vartheta_\sigma*\widetilde{\Phi}_\varepsilon.
\]
For $0<\sigma<\delta_0/4$, the restriction of $\Psi_{\varepsilon,\sigma}$ to
$\mathbb R^3\times[0,T_0)$ belongs to
$C_c^\infty(\mathbb R^3\times[0,T_0))$.

Moreover, as $\sigma\to0$,
\[
\Psi_{\varepsilon,\sigma}\to \Phi_\varepsilon
\qquad\text{uniformly on }\mathbb R^3\times[0,T_0],
\]
\[
\partial_t\Psi_{\varepsilon,\sigma}\to \partial_t\Phi_\varepsilon,
\qquad
\nabla\Psi_{\varepsilon,\sigma}\to \nabla\Phi_\varepsilon
\quad\text{a.e. on }\mathbb R^3\times(0,T_0),
\]
and
\[
\Psi_{\varepsilon,\sigma}(\cdot,0)\to \Phi_\varepsilon(\cdot,0)
\qquad\text{uniformly on }\mathbb R^3.
\]
These functions are uniformly bounded, and their supports are contained in one fixed compact
subset of $\mathbb R^3\times[0,T_0)$; likewise,
$\Psi_{\varepsilon,\sigma}(\cdot,0)$ are supported in one fixed compact subset of
$\mathbb R^3$.

Testing \eqref{eq:transport_distributional_appendix} with
$\Psi_{\varepsilon,\sigma}$ and passing to the limit $\sigma\to0$ by dominated convergence,
using $\xi,u\xi\in L^1_{\mathrm{loc}}$ and $\xi_0\in L^1$, we obtain
\[
0=
\int_0^{T_0}\int_{\mathbb R^3}
\xi\bigl(\partial_t\Phi_\varepsilon+u\cdot\nabla\Phi_\varepsilon\bigr)\,dx\,dt
+\int_{\mathbb R^3}\xi_0(x)\Phi_\varepsilon(x,0)\,dx .
\]
Using \eqref{eq:Phi_eps_eq_appendix},
\begin{equation}
\label{eq:key_identity_appendix}
\int_0^{T_0}\int_{\mathbb R^3}\xi(x,t)\chi(t)\eta_\varepsilon(x)\,dx\,dt
=
\int_0^{T_0}\int_{\mathbb R^3}\xi(x,t)R_\varepsilon(x,t)\,dx\,dt
+\int_{\mathbb R^3}\xi_0(x)\Phi_\varepsilon(x,0)\,dx .
\end{equation}


\smallskip
\noindent\emph{Step 5: passage to the limit.}
Because $\eta_\varepsilon\to\eta$ uniformly and $\operatorname{supp}\eta_\varepsilon$ stays
inside a fixed compact set, while $\xi\in L^1_{\mathrm{loc}}$, the left-hand side of
\eqref{eq:key_identity_appendix} converges to
\[
\int_0^{T_0}\int_{\mathbb R^3}\xi(x,t)\chi(t)\eta(x)\,dx\,dt.
\]
Also, by \eqref{eq:R_eps_support_appendix}, for $0<\varepsilon<1$ we have
\[
\supp R_\varepsilon(\cdot,t)\subset K_*+\overline B_1
\qquad\text{for every }t\in[0,T_0].
\]
Hence, using \eqref{eq:R_eps_to_zero_appendix} and $\xi\in L^1_{\mathrm{loc}}$, we obtain
\[
\left|
\int_0^{T_0}\int_{\mathbb R^3}\xi(x,t)R_\varepsilon(x,t)\,dx\,dt
\right|
\le
\|R_\varepsilon\|_{L^\infty((0,T_0)\times\mathbb R^3)}
\int_0^{T_0}\int_{K_*+\overline B_1}|\xi(x,t)|\,dx\,dt
\longrightarrow 0.
\]


Finally, since $\Phi(\cdot,0)$ is compactly supported and H\"older continuous,
$\Phi_\varepsilon(\cdot,0)\to\Phi(\cdot,0)$ uniformly, hence
\begin{align*}
\int_{\mathbb R^3}\xi_0(x)\Phi_\varepsilon(x,0)\,dx
&\longrightarrow
\int_{\mathbb R^3}\xi_0(a)\Phi(a,0)\,da \\
&=
\int_{\mathbb R^3}\xi_0(a)\int_0^{T_0}\chi(\tau)\eta\big(Y_\tau(a)\big)\,d\tau\,da \\
&=
\int_0^{T_0}\chi(\tau)\int_{\mathbb R^3}\xi_0(a)\eta\big(Y_\tau(a)\big)\,da\,d\tau 
=
\int_0^{T_0}\chi(\tau)\int_{\mathbb R^3}\widetilde{\xi}(x,\tau)\eta(x)\,dx\,d\tau,
\end{align*}
where in the last step we used the change of variables $x=Y_\tau(a)$ and the
measure-preserving property of $Y_\tau$.

Passing to the limit $\varepsilon\to0$ in \eqref{eq:key_identity_appendix}, we obtain
\[
\int_0^{T_0}\int_{\mathbb R^3}\xi(x,t)\chi(t)\eta(x)\,dx\,dt
=
\int_0^{T_0}\int_{\mathbb R^3}\widetilde{\xi}(x,t)\chi(t)\eta(x)\,dx\,dt.
\]
Since $\chi\in C_c^\infty(0,T_0)$ and $\eta\in C_c^\infty(\mathbb R^3)$ are arbitrary,
it follows that
\[
\xi=\widetilde{\xi}
\qquad\text{in }\mathcal D'((0,T_0)\times\mathbb R^3),
\]
hence a.e. on $(0,T_0)\times\mathbb R^3$ because both functions are locally integrable.
This proves
\[
\xi(x,t)=\xi_0\bigl(Y_t^{-1}(x)\bigr)
\qquad\text{for a.e. }(x,t)\in\mathbb R^3\times(0,T_0),
\]
and completes the proof.
\end{proof}

\end{document}
